\documentclass[10pt,a4paper]{amsart}
\usepackage[margin=19mm]{geometry}
\usepackage{amsmath,amssymb,mathtools}
\usepackage[scr=rsfs,cal=euler]{mathalfa}
\usepackage{xfrac}
\usepackage[unicode,hidelinks,bookmarks=false]{hyperref}
\newcommand{\ie}{i.e.\ }
\newcommand{\cf}{cf.\ }
\newcommand{\resp}{respectively\ }
\newcommand{\Subsection}{\subsection}
\providecommand{\nobreakdash}{\nobreak\hbox{-}}
\setlength{\emergencystretch}{2em}
\multlinegap0pt
\usepackage{amssymb}
\usepackage{enumerate}
\usepackage{xcolor}
\usepackage{tikz}
\newtheorem{defn}{Definition}
\newtheorem{theorem}{Theorem}
\newcommand{\R}{{ \mathbb{R}  }}
\newcommand{\abs}[1]{\left| #1 \right|}
\newcommand{\bke}[1]{\left( #1 \right)}
\newcommand{\norm}[1]{\left\Vert #1 \right\Vert}
\title{Existence of weak solutions for fast diffusion equation \\ with a divergence type of drift term}
\author{Sukjung Hwang, Kyungkeun Kang, Hwa Kil Kim}
\date{Source: arXiv:2501.09539v3 (CC BY 4.0)}
\begin{document}
\maketitle

\noindent\emph{Source and licence.} Existence of weak solutions for fast diffusion equation \\ with a divergence type of drift term. Version arXiv:2501.09539v3 (CC BY 4.0), distributed by arXiv under the Creative Commons Attribution 4.0 International licence, http://creativecommons.org/licenses/by/4.0/, as recorded in the arXiv licence metadata for this exact version and shown on the abstract page https://arxiv.org/abs/2501.09539v3. Full licence notice: \texttt{licenses/arxiv-2501.09539-CC-BY-4.0.txt}.\par\smallskip

\noindent\textit{Editorial scope.} Counted unit: the article's theorem BS-energy (source0.tex lines 1185--1193). The article's complete original proof of it is delivered with it (source0.tex lines 1194--1255, one \texttt{proof} environment of 62 lines). No mathematics is written, repaired or re-derived: the statement and the proof are the article's own lines, in the article's own notation, and no retained line is substituted or re-flowed.  Retained uncounted as the source-local context the proof needs: lines 1128--1130; lines 1131--1133; lines 1150--1165.  Nothing was omitted from any retained span, so the package carries no omission record.  No retained line contains a citation, so the wrapper carries no bibliography and no bibliography entry.  No retained line contains a figure environment, an image-inclusion command or drawing code, so no image is delivered and the package declares no asset dependency.  Editorial formatting only, in addition: an AMS \texttt{amsart} preamble that restates the article's own declarations and macros used by the retained lines, namely the environments \texttt{defn}, \texttt{theorem} on separate counters, together with the standard packages those lines need.  The retained environments print in this document as Definition 1, Theorem 1 in order of appearance, whereas the article prints the same items with the numbering of its own full document.  The complete native source of the article as distributed in the arXiv e-print for this exact version is bundled unchanged as \texttt{sources/171685/source0.tex}.
\begin{equation}\label{BE-10}
\theta_t-\Delta \theta^m +u\cdot \nabla \theta=0,
\end{equation}

\begin{equation}\label{BE-20}
u_t-\Delta u+(u\cdot \nabla) u+\nabla \pi=-\theta {\bf{e}_d},\qquad {\rm div}\,u=0.
\end{equation}

\begin{defn}\label{Weak-ks} Let $0<m<1$ and $\Omega\subset\R^d$, $d=2,3$ be a bounded domain with smooth boundaries. We say that a pair of $(\theta, u)$ is a  weak solution of \eqref{BE-10}-\eqref{BE-20} in $\Omega_{T} := \Omega \times (0, T)$ with zero boundary data if the followings are satisfied:
\begin{itemize}
	\item [(i)] It holds that 
	\[
	\theta , \ \nabla \theta^m, \ \theta u , \ u\otimes u, \ \nabla u \in L_{x,t}^{1}.
	\]
	
	\item [(ii)] For any $\varphi \in \mathcal{C}_{c}^{\infty} \left(\overline{\Omega} \times [0, T);\R\right)$ and $\psi \in \mathcal{C}_{c}^{\infty} \left(\overline{\Omega} \times [0, T);\R^d\right)$ with $\rm{div}\,\psi=0$ both of which vanish on $\partial\Omega$, it holds that 
	\[
	\iint_{\Omega_{T}} \left\{ - \theta \varphi_t + \nabla \theta^m \nabla \varphi - \theta u\cdot \nabla \varphi \right\} \,dxdt = \int_{\Omega}\theta_0 \varphi(\cdot,0) \,dx,
	\]
	\[
	\iint_{\Omega_{T}} \left\{ - u \psi_t + \nabla u \nabla \psi - u\otimes u:\nabla \psi \right\} \,dxdt = \int_{\Omega}u_0 \psi(\cdot,0) \,dx-\iint_{\Omega_{T}}  \theta {\bf e}_d  \psi \,dxdt,
	\]
\end{itemize}	
\end{defn}

\begin{theorem}
Let  $\Omega\subset\R^d$, $d=2,3$ be a bounded domain with smooth boundaries. 
If $0< m <1$ when $d=2$ and if $\frac{2}{3} \le m<1$ when $d=3$, then,
there exists a pair of weak solution $(\theta, u)$ in Definition \ref{Weak-ks}. Furthermore, for any  $(\theta, u)$ satisfies 
\begin{equation}\label{BS-energy}
\sup_{t\in (0, T)} \int_{\Omega} (\theta(\cdot, t) +\abs{u(\cdot, t)}^2) \,dx
+\iint_{\Omega_T} (\left|\nabla \theta^{\frac{m}{2}}\right|^{2} +\abs{\nabla u}^2) \,dxdt\le C\bke{\int_{\Omega}\theta_0 \log \theta_0 \,dx,  \norm{\theta_0}_{L^1}, \norm{u_0}_{L^{2}_{x}}}.
\end{equation}
\end{theorem}

\begin{proof}
We first note that 
\begin{equation}\label{ks-pf-100}
\frac{d}{dt}\int_{\Omega} \theta(\cdot, t)\log \theta(\cdot, t) \,dx+\int_{\Omega} \left|\nabla \theta^{\frac{m}{2}}\right|^{2} \,dx= 0
\end{equation}

We treat the two and three dimensions separately.
\\
$\bullet$ \underline{(2d case)}\,\,
Let $p\in (1, \infty)$ with $\frac{2}{m}\le p$.
For the fluid equation, energy estimate gives
\[
\frac{1}{2}\frac{d}{dt}\norm{u}_{L^{2}_{x}}^2+\norm{\nabla u}_{L^{2}_{x}}^2\lesssim \int_{\Omega}\theta \abs{u}dx\lesssim \norm{\theta}_{L^{q}_{x}}\norm{u}_{L^{p}_{x}} \lesssim \norm{\theta}_{L^{q}_{x}}\norm{u}^{\frac{2}{p}}_{L^{2}_{x}} \norm{\nabla u}^{\frac{p-2}{p}}_{L^{2}_{x}},
\]
where $\frac{1}{q}+ \frac{1}{p}=1$.
Due to Young's inequality, we have
\[
\frac{d}{dt}\norm{u}_{L^{2}_{x}}^2+\norm{\nabla u}_{L^{2}_{x}}^2\lesssim \norm{\theta}^{\frac{2p}{p+2}}_{L^{q}_{x}}\norm{u}^{\frac{4}{p+2}}_{L^{2}_{x}} =\norm{\theta^{\frac{m}{2}}}^{\frac{4p}{m(p+2)}}_{L^{\frac{2q}{m}}_{x}}\norm{u}^{\frac{4}{p+2}}_{L^{2}_{x}}.
\]
Noting that $\norm{\theta(t)}_{L^1}=\norm{\theta_0}_{L^1}$ and
\[
\norm{\theta^{\frac{m}{2}}}^{\frac{4p}{m(p+2)}}_{L^{\frac{2q}{m}}_{x}}\le \bke{\norm{\theta^{\frac{m}{2}}}^{\frac{1}{q}}_{L^{\frac{2}{m}}_{x}}\norm{\nabla \theta^{\frac{m}{2}}}^{\frac{1}{p}}_{L^{2}_{x}}}^{\frac{4p}{m(p+2)}}\lesssim\norm{\nabla \theta^{\frac{m}{2}}}^{\frac{4}{m(p+2)}}_{L^{2}_{x}},
\]
we obtain
\begin{equation}\label{ks-pf-110}
\frac{d}{dt}\norm{u}_{L^{2}_{x}}^2+\norm{\nabla u}_{L^{2}_{x}}^2\lesssim \norm{\nabla \theta^{\frac{m}{2}}}^{\frac{4}{m(p+2)}}_{L^{2}_{x}}\norm{u}^{\frac{4}{p+2}}_{L^{2}_{x}}
\lesssim \epsilon\norm{\nabla \theta^{\frac{m}{2}}}^{2}_{L^{2}_{x}}+\norm{u}^{\frac{4m}{m(p+2)-2}}_{L^{2}_{x}}.
\end{equation}
Combining \eqref{ks-pf-100} and \eqref{ks-pf-110}, we obtain that 
\begin{equation*}
\frac{d}{dt}\int_{\Omega} (\theta \log \theta +\abs{u}^2)\,dx+\iint_{\Omega_T} (\left|\nabla \theta^{\frac{m}{2}}\right|^{2} +\abs{\nabla u}^2) \,dxdt\lesssim \norm{u}^{\frac{4m}{m(p+2)-2}}_{L^{2}_{x}}.
\end{equation*}
It is required that 
\[
\frac{4m}{m(p+2)-2}\le 2 \quad \Longrightarrow \quad \frac{2}{m} \le p.
\]
Since $u\in L^{a, b}_{x,t}$ where $\frac{2}{a}+\frac{2}{b}=1$ with $2\le a<\infty$, it is automatic that $u\in\mathcal{D}_{m,1}^{(q_1, q_2)}$ with $d=2$.
Indeed, in case $m\le \frac{1}{2}$, it is straightforward. If $m \in (0, \frac{1}{2})$, then we can take $q_2 \le \frac{2(1-m)}{(1-2m)}$ and then choose $q_1$ with $\frac{2}{q_1}+\frac{2}{q_2}=1$, which implies that $u\in\mathcal{D}_{m,1}^{(q_1, q_2)}$, $d=2$.
This completes the proof for 2d case.
\\
\\
$\bullet$ \underline{(3d case)}\,\, 
%Let $m>\frac{2}{3}$. 
As similarly, we have 
\begin{equation}\label{ks-pf-120}
\frac{1}{2}\frac{d}{dt}\norm{u}_{L^{2}_{x}}^2+\norm{\nabla u}_{L^{2}_{x}}^2\lesssim \int_{\Omega}\theta \abs{u}dx\lesssim \norm{\theta}_{L^{\frac{6}{5}}_{x}}\norm{u}_{L^{6}_{x}} \lesssim \norm{\theta}^{2}_{L^{\frac{6}{5}}_{x}} +\frac{1}{2}\norm{\nabla u}^2_{L^{2}_{x}}
\end{equation}
Via interpolation, it gives
\begin{equation}\label{ks-pf-130}
\norm{\theta}^{2}_{L^{\frac{6}{5}}_{x}}=\norm{\theta}^{\frac{4}{m}}_{L^{\frac{12}{5m}}_{x}}\le \norm{\theta^{\frac{m}{2}}}^{\frac{10m-4}{m(3m-1)}}_{L^{\frac{2}{m}}_{x}}\norm{\nabla \theta^{\frac{m}{2}}}^{\frac{2}{3m-1}}_{L^{2}_{x}}\lesssim \norm{\nabla \theta^{\frac{m}{2}}}^{\frac{2}{3m-1}}_{L^{2}_{x}}.
\end{equation}
%As assumed at the beginning, this implies $m>\frac{2}{3}$.
It is required that
\[
\frac{2}{3m-1}\le 2 \quad \Longrightarrow \quad  m \geq \frac{2}{3}.
\]
We note that $u\in L^{6,2}_{x,t}$, which belongs to $u\in\mathcal{D}_{m,1}^{(q_1, q_2)}$, $d=3$. Indeed, since $u\in L^{6,2}_{x,t}$, it is direct that 
\[
\frac{3}{6} + \frac{2+3(m-1)}{2} \leq 2+ 3(m-1)\quad\Longrightarrow\quad \frac{2}{3}\leq m.
\]
Combining \eqref{ks-pf-100}, \eqref{ks-pf-120} and \eqref{ks-pf-130}, we complete the proof.
\end{proof}

\end{document}
